\section{Computational results}

The following statements are exhaustive finite computations.  Their inputs,
programs, commands, and outputs are specified in
Section~\ref{sec:reproducibility}.

\begin{computationaltheorem}[Small orders]\label{cthm:small-orders}
The unique reduced Latin square of order $2$ is FFF, and all four reduced
Latin squares of order $4$ are FFF.  None of the $9408$ reduced Latin squares
of order $6$ is FFF.  More precisely, their patterns are
\[
 6600\,\TTT,\qquad 936\,\mathsf{TFF},\qquad
 936\,\mathsf{FTF},\qquad 936\,\mathsf{FFT}.
\]
\end{computationaltheorem}

The generator constructs reduced Latin squares by exact cover and loads no
cached square list.  Every generated square was checked to be Latin and
reduced, all $9408$ table strings were distinct, and the generated set agrees
exactly with the ANU reduced-order-$6$ data file~\cite{McKayLatinData}.
Proposition~\ref{prop:invariance} therefore extends the nonexistence statement
from reduced squares to all Latin squares of order $6$.

For all order-$8$ record references below, ``local source line'' means the
one-based physical line number in the frozen file consumed by the scanner.
The ANU catalogue numbers its records from zero, so local source line $j$
corresponds to ANU record $j-1$.

\begin{computationaltheorem}[Order-$8$ pattern census]\label{cthm:order8}
Across the $283657$ order-$8$ main-class representatives, the pattern
distribution is
\begin{center}
\begin{tabular}{cr@{\qquad}cr@{\qquad}cr@{\qquad}cr}
\toprule
Pattern & Count & Pattern & Count & Pattern & Count & Pattern & Count\\
\midrule
$\FFF$ & 230 & $\mathsf{FFT}$ & 81 & $\mathsf{FTF}$ & 40 &
$\mathsf{FTT}$ & 315\\
$\mathsf{TFF}$ & 133 & $\mathsf{TFT}$ & 209 & $\mathsf{TTF}$ & 232 &
$\TTT$ & 282417\\
\bottomrule
\end{tabular}
\end{center}
Among the $230$ FFF main classes, exactly $5$ are group-isotopic and $225$
are not group-isotopic.
\end{computationaltheorem}

\begin{computationaltheorem}[Subsquare census of the FFF classes]
Among the $230$ FFF representatives, the number of intercalates ranges from
$2$ to $112$.  Each endpoint is attained by exactly one representative: local
source line $10776$ (ANU record $10775$) is the unique minimum class and is not
group-isotopic, while local source line $1$ (ANU record $0$) is the unique
maximum class and is group-isotopic.  No representative
contains an order-$3$ subsquare.  The distribution of the number of
order-$4$ subsquares is
\[
  0^{63},\qquad 4^{156},\qquad 12^{10},\qquad 28^1,
\]
where exponents count representatives.
\end{computationaltheorem}

The cycle-parity obstruction of Proposition~\ref{prop:q-obstruction} is also
scanned on the complete order-$8$ input.  Exactly $1285$ representatives are
QQQ: the $230$ FFF representatives and $1055$ non-FFF representatives.  This
is an exact census fact and demonstrates that QQQ is necessary but not
sufficient for FFF.

\begin{computationaltheorem}[Affine cycle-profile saturation]
\label{cthm:affine-profile-saturation}
Over each of $\operatorname{GF}(2)$, $\operatorname{GF}(3)$,
$\operatorname{GF}(5)$, and $\operatorname{GF}(7)$, every affine relation
among the aggregate cycle profiles of all $283657$ order-$8$ main-class
representatives lies in the span of the relations in
Proposition~\ref{prop:aggregate-profile-baseline}.
\end{computationaltheorem}

Incremental finite-field row reduction includes a constant coordinate and
processes the complete census directly.  The analogous order-$6$ audit
matches the baseline over the three odd fields but has three additional
Latin-specific relations over $\operatorname{GF}(2)$.  Conversely, three
extra binary relations visible in the tracked $135$-table order-$10$ partial
corpus are refuted by explicit Latin squares obtained from a deterministic
intercalate-trade walk on the $D_{10}$ table.  Thus those order-$10$
relations were sampling artifacts.  This does not constitute an exhaustive
order-$10$ classification.

The associated validator checks every order-$2$, order-$3$, and order-$4$
subsquare and finds no violation of
Theorem~\ref{thm:subsquare-monotonicity}.  The multiples of four in the
order-$4$ distribution also follow directly: an order-$m$ subsquare inside
an order-$2m$ square forces the three complementary blocks to be Latin
subsquares, so half-order subsquares occur in complementary quartets.

The non-symmetric labels in the table refer to the particular representatives
in the cited ANU file.  Parastrophy permutes the three bits, so, for example,
an ordered label such as $\mathsf{FFT}$ is not itself a main-class invariant.
The FFF count, the existence of every pattern, and the group-isotopy split are
unaffected by this distinction.

The scan tests every unordered pair of lines in all three views.  A separate
closure-family computation reproduces the $5/225$ split using
Proposition~\ref{prop:closure}.  The input is the ANU/McKay main-class file,
whose stated size and equivalence convention are given on the data
page~\cite{McKayLatinData}; our frozen copy is an exact byte match to the file
served by that URL on August 9, 2026.

\begin{corollary}\label{cor:all-patterns}
For every $m\geq3$, Latin squares of order $2^m$ realize all eight patterns.
Every such order also contains a non-group-isotopic FFF Latin square.
\end{corollary}

\begin{proof}
Computational Theorem~\ref{cthm:order8} supplies one order-$8$ square of each
pattern.  Take direct products with the FFF group table of $C_2^{m-3}$.
The pattern statement follows from Theorem~\ref{thm:pattern-product}.  For
the second assertion, instead start from the explicit loop in
Theorem~\ref{thm:explicit-nongroup-loop}; non-group-isotopicity is preserved
by Theorem~\ref{thm:group-product}.
\end{proof}

This is an existence result only.  It neither counts higher-order main
classes nor asserts that direct products of distinct order-$8$ classes remain
distinct.

As an additional check, Wanless' complete online list of order-$8$ species
with transitive autotopism group contains $11$ squares~\cite{WanlessTransitiveData}.
The frozen local input is a byte-for-byte match to the online file (SHA-256
\nolinkurl{1c777328c099e3c242a49f9f2019d83314d7ab800505c7d9291bdb16ef029146}).
All $11$ have pattern FFF; $5$ are group-isotopic and $6$ are not.  This
secondary computation is not needed for Corollary~\ref{cor:all-patterns}.
